\documentclass[10pt]{article}

\usepackage{graphics}
\usepackage{hevea}
\usepackage{verbatim}


\newcommand{\imgscale}{0.7}


\newcommand{\ns}[1]{\texttt{#1}}
\newcommand{\ejabberd}{\texttt{ejabberd}}
\newcommand{\Jabber}{Jabber}

\newcommand{\modregister}{\texttt{mod\_register}}
\newcommand{\modroster}{\texttt{mod\_roster}}
\newcommand{\modconfigure}{\texttt{mod\_configure}}
\newcommand{\moddisco}{\texttt{mod\_disco}}
\newcommand{\modstats}{\texttt{mod\_stats}}
\newcommand{\modvcard}{\texttt{mod\_vcard}}
\newcommand{\modoffline}{\texttt{mod\_offline}}
\newcommand{\modecho}{\texttt{mod\_echo}}
\newcommand{\modprivate}{\texttt{mod\_private}}
\newcommand{\modtime}{\texttt{mod\_time}}
\newcommand{\modversion}{\texttt{mod\_version}}



\title{Ejabberd Installation and Operation Guide}
\author{Alexey Shchepin \\
  \ahrefurl{mailto:alexey@sevcom.net} \\
  \ahrefurl{xmpp:aleksey@jabber.ru}}
\date{January 26, 2003}

\begin{document}
\begin{titlepage}
  \maketitle{}
  
  {\centering
    \imgsrc{logo.png}{}
    \begin{latexonly}
      \scalebox{\imgscale}{\includegraphics{logo.png}}
    \end{latexonly}
    \par
  }
\end{titlepage}
%\newpage
\tableofcontents{}

\newpage
\section{Introduction}
\label{sec:intro}

\ejabberd{} is a Free and Open Source distributed fault-tolerant \Jabber{}
server.  It writen mostly in Erlang.

TBD



\section{Installation}
\label{sec:installation}


\subsection{Installation Requirements}
\label{sec:installreq}

To compile \ejabberd{}, you need following packages:
\begin{itemize}
\item GNU Make;
\item GCC;
\item libexpat 1.95 or later;
\item Erlang/OTP R8B or later.
\end{itemize}

\subsection{Obtaining}
\label{sec:obtaining}

Currently no stable version released.

Latest alpha version can be retrieved via CVS\@.  Do following steps:
\begin{itemize}
\item \texttt{export CVSROOT=:pserver:cvs@www.jabber.ru:/var/spool/cvs}
\item \texttt{cvs login}
\item Enter empty password
\item \texttt{cvs -z3 co ejabberd}
\end{itemize}






\subsection{Compilation}
\label{sec:compilation}




%\subsection{Initial Configuration}
%\label{sec:initconfig}


\section{Configuration}
\label{sec:configuration}

\subsection{Initial Configuration}
\label{sec:initconfig}

%\verbatiminput{../src/ejabberd.cfg}

Configuration file is loaded after first start of \ejabberd{}.  It consists of
sequence of Erlang terms.  Parts of lines after \texttt{`\%'} sign are ignored.
Each term is tuple, where first element is name of option, and other are option
values.


\subsubsection{Host Name}
\label{sec:confighostname}

Option \texttt{hostname} defines name of \Jabber{} domain that \ejabberd{}
serves.  E.\,g.\ to use \texttt{jabber.org} domain add following line in config:
\begin{verbatim}
{host, "jabber.org"}.
\end{verbatim}

%This option is mandatory.



\subsubsection{Listened Sockets}
\label{sec:configlistened}

Option \texttt{listen} defines list of listened sockets and what services
runned on them.  Each element of list is a tuple with following elements:
\begin{itemize}
\item Port number;
\item Module that serves this port;
\item Function in this module that starts connection (likely will be removed);
\item Options to this module.
\end{itemize}

Currently three modules implemented:
\begin{itemize}
\item \texttt{ejabberd\_c2s}: serves C2S connections;
\item \texttt{ejabberd\_s2s\_in}: serves incoming S2S connections;
\item \texttt{ejabberd\_service}: serves connections to \Jabber{} services
  (i.\,e.\ that use \texttt{jabber:component:accept} namespace).
\end{itemize}

For example, following configuration defines that C2S connections listened on
port 5222, S2S on port 5269 and that service \texttt{conference.jabber.org}
must be connected to port 8888 with password ``\texttt{secret}''.

\begin{verbatim}
{listen, [{5222, ejabberd_c2s,     start, []},
          {5269, ejabberd_s2s_in,  start, []},
          {8888, ejabberd_service, start, ["conference.jabber.org", "secret"]}
         ]}.
\end{verbatim}


\subsubsection{Access Rules}
\label{sec:configaccess}

Access control in \ejabberd{} is done via Access Control Lists (ACL).  In
config file they looks like this:
\begin{verbatim}
{acl, <aclname>, {<acltype>, ...}}.
\end{verbatim}

\texttt{<acltype>} can be one of following:
\begin{description}
\item[\texttt{all}] Matches all JIDs.  Example:
\begin{verbatim}
{acl, all, all}.
\end{verbatim}
\item[\texttt{\{user, <username>\}}] Matches local user with name
  \texttt{<username>}.  Example:
\begin{verbatim}
{acl, admin, {user, "aleksey"}}.
\end{verbatim}
\item[\texttt{\{user, <username>, <server>\}}] Matches user with JID
  \texttt{<username>@<server>}.  Example:
\begin{verbatim}
{acl, admin, {user, "aleksey", "jabber.ru"}}.
\end{verbatim}
\item[\texttt{\{server, <server>\}}] Matches any JID from server
  \texttt{<server>}.  Example:
\begin{verbatim}
{acl, jabberorg, {server, "jabber.org"}}.
\end{verbatim}
\end{description}

Allowing or denying of different services is like this:
\begin{verbatim}
{access, <accessname>, [{allow, <aclname>},
                        {deny, <aclname>},
                        ...
                       ]}.
\end{verbatim}
When JID is checked to have access to \texttt{<accessname>}, server
sequentially checks if this JID in one of the ACLs that are second elements in
eache tuple in list.  If one of them matched, then returned first element of
matched tuple.  Else returned ``\texttt{deny}''.

Example:
\begin{verbatim}
{access, configure, [{allow, admin}]}.
{access, something, [{deny, badmans},
                     {allow, all}]}.
\end{verbatim}



\subsubsection{Modules}
\label{sec:configmodules}

Option \texttt{modules} defines list of modules that will be loaded after
\ejabberd{} startup.  Each list element is a tuple where first element is a
name of module and second is list of options to this module.  See
section~\ref{sec:modules} for detailed information on each module.

Example:
\begin{verbatim}
{modules, [
           {mod_register,  []},
           {mod_roster,    []},
           {mod_configure, []},
           {mod_disco,     []},
           {mod_stats,     []},
           {mod_vcard,     []},
           {mod_offline,   []},
           {mod_echo,      [{host, "echo.localhost"}]},
           {mod_private,   []},
           {mod_time,      [{iqdisc, no_queue}]},
           {mod_version,   []}
          ]}.
\end{verbatim}


\subsection{Online Configuration}
\label{sec:onlineconfig}

To use facility of online reconfiguration of \ejabberd{} needed to have
\modconfigure{} loaded (see section~\ref{sec:modconfigure}).  Also highly
recommended to load \moddisco{} (see section~\ref{sec:moddisco}), because
\modconfigure{} highly integrates with it.  Also recommended to use
disco- and xdata-capable client.

TBD

\section{Distribution}
\label{sec:distribution}


\subsection{How it works}
\label{sec:howitworks}

\Jabber{} domain is served by one or more \ejabberd{} nodes.  This nodes can be
runned on different machines that can be connected via network.  They all must
have access to connect to port 4369 of all another nodes, and must have same
magic cookie (see Erlang/OTP documentation, in short file
\texttt{\~ejabberd/.erlang.cookie} must be the same on all nodes).  This is
needed because all nodes exchange information about connected users, S2S
connection, registered services, etc\ldots

Each \ejabberd{} node run following modules:
\begin{itemize}
\item router;
\item local router.
\item session manager;
\item S2S manager;
\end{itemize}


\subsubsection{Router}

This module is the main router of \Jabber{} packets on each node.  It route
them based on their destanations domains.  It have two tables: local and global
routes.  First, domain of packet destination searched in local table, and if it
finded, then packet routed to appropriate process.  If no, then it searched in
global table, and routed to appropriate \ejabberd{} node or process.  If it not
exists in both tables, then it sended to S2S manager.


\subsubsection{Local Router}

This module route packets which have destination domain equal to this server
name.  If destination JID have node, then it routed to session manager, else it
processed depending on it content.


\subsubsection{Session Manager}

This module route packets to local users.  It search to what user resource
packet must be sended via presence table.  If this reseouce connected to this
node, it routed to C2S process, if it connected via another node, then packet
sended to session manager on it.


\subsubsection{S2S Manager}

This module route packets to another \Jabber{} servers.  First, it check if
to domain of packet destination from domain of source already opened S2S
connection.  If it opened on another node, then it routed to S2S manager on
that node, if it opened on this node, then it routed to process that serve this
connection, and if this connection not exists, then it opened and registered.


\section{Built-in Modules}
\label{sec:modules}



\subsection{Common Options}
\label{sec:modcommonopts}

Following options used by many modules, so they described in separate section.


\subsubsection{Option \texttt{iqdisc}}

Many modules define handlers for processing IQ queries of different namespaces
to this server or to user (e.\,g.\ to \texttt{myjabber.org} or to
\texttt{user@myjabber.org}).  This option defines processing discipline of this
queries.  Possible values are:
\begin{description}
\item[\texttt{no\_queue}] All queries of namespace with this processing
  discipline processed immediately.  This also means that no other packets can
  be processed until finished this.  Hence this discipline is not recommended
  if processing of query can take relative many time.
\item[\texttt{one\_queue}] In this case created separate queue for processing
  IQ queries of namespace with this discipline, and processing of this queue
  done in parallel with processing of other packets. This discipline is most
  recommended.
\item[\texttt{parallel}] In this case for all packets of namespace with this
  discipline spawned separate Erlang process, so all this packets processed in
  parallel.  Although spawning of Erlang process have relative low cost, this
  can broke server normal work, because Erlang have limit of 32000 processes.
\end{description}

Example:
\begin{verbatim}
{modules, [
           ...
           {mod_time,      [{iqdisc, no_queue}]},
           ...
          ]}.
\end{verbatim}


\subsubsection{Option \texttt{host}}

Some modules may act as services, and wants to have different domain name.
This option explicitly defines this name.

Example:
\begin{verbatim}
{modules, [
           ...
           {mod_echo,      [{host, "echo.myjabber.org"}]},
           ...
          ]}.
\end{verbatim}



\subsection{\modregister{}}
\label{sec:modregister}



\subsection{\modroster{}}
\label{sec:modroster}



\subsection{\modconfigure{}}
\label{sec:modconfigure}



\subsection{\moddisco{}}
\label{sec:moddisco}



\subsection{\modstats{}}
\label{sec:modstats}

This module adds support of
\footahref{http://www.jabber.org/jeps/jep-0039.html}{JEP-0039} (Statistics Gathering).

Options:
\begin{description}
\item[\texttt{iqdisc}] \ns{http://jabber.org/protocol/stats} IQ queries
  processing discipline.
\end{description}

TBD about access.

\subsection{\modvcard{}}
\label{sec:modvcard}



\subsection{\modoffline{}}
\label{sec:modoffline}



\subsection{\modecho{}}
\label{sec:modecho}



\subsection{\modprivate{}}
\label{sec:modprivate}

This module adds support of
\footahref{http://www.jabber.org/jeps/jep-0049.html}{JEP-0049} (Private XML
Storage).

Options:
\begin{description}
\item[\texttt{iqdisc}] \ns{jabber:iq:private} IQ queries processing discipline.
\end{description}

\subsection{\modtime{}}
\label{sec:modtime}

This module answers UTC time on \ns{jabber:iq:time} queries.

Options:
\begin{description}
\item[\texttt{iqdisc}] \ns{jabber:iq:time} IQ queries processing discipline.
\end{description}


\subsection{\modversion{}}
\label{sec:modversion}

This module answers \ejabberd{} version on \ns{jabber:iq:version} queries.

Options:
\begin{description}
\item[\texttt{iqdisc}] \ns{jabber:iq:version} IQ queries processing discipline.
\end{description}






\end{document}
